%% ========================================================================
%% chapter05-process-management.tex
%% Manajemen Proses di Linux
%% ========================================================================

\chapter{Manajemen Proses}
\label{ch:process-management}

\chapterhero{red}{Amati bagaimana proses dibuat, dijadwalkan, diprioritaskan, dan dihentikan agar kamu bisa mengendalikan beban kerja sistem dengan tepat.}{\faTasks\quad process}{\faBell\quad signal}{\faSortAmountDown\quad priority}

Saat kamu membuka aplikasi musik di ponsel sambil membaca pesan dan men-download foto,
ponsel itu menjalankan beberapa program sekaligus. Tidak ada yang terasa lambat karena
sistem operasi mengatur semua program tersebut dengan rapi: memberi giliran waktu prosesor,
mengalokasikan memori, dan memastikan setiap program mendapat sumber daya yang dibutuhkan.
Linux bekerja persis seperti itu, hanya saja seluruh mekanismenya bisa kamu lihat dan
kendalikan langsung dari terminal.

Di Linux, setiap program yang sedang berjalan disebut \textbf{proses}. Ketika kamu
mengetik \cmd{ls} di terminal, Linux membuat proses baru, menjalankan program tersebut,
lalu proses itu selesai dan menghilang. Setiap proses diberi nomor unik yang disebut
\textbf{PID} (\textit{Process ID}). Proses yang membuat proses lain disebut proses induk,
dan proses yang dibuat disebut proses anak. Seluruh proses di sistem membentuk hierarki
seperti pohon keluarga, dengan proses pertama yang dijalankan kernel saat booting sebagai
leluhur semua proses lainnya.

Memahami manajemen proses bukan hanya pengetahuan teoritis. Ketika sebuah aplikasi tiba-tiba
memakan hampir semua kapasitas prosesor, kamu perlu tahu cara menemukannya dan menghentikannya
tanpa mematikan sistem. Ketika proses tertentu harus berjalan dengan prioritas lebih rendah
agar tidak mengganggu layanan utama, ada perintah khusus untuk mengaturnya. Ketika kamu
perlu menjalankan sesuatu di latar belakang dan menutup terminal tanpa menghentikan prosesnya,
ada cara yang tepat untuk itu. Semua kemampuan ini akan kamu kuasai di bab ini.

\textbf{Yang Akan Kamu Pelajari}

Setelah menyelesaikan bab ini, kamu akan mampu:

\begin{enumerate}
    \item menjelaskan konsep proses dan thread serta perbedaan keduanya di Linux;
    \item membaca status proses dan memahami siklus hidup proses dari lahir hingga selesai;
    \item mengatur prioritas proses menggunakan \cmd{nice} dan \cmd{renice};
    \item mengirim sinyal ke proses menggunakan \cmd{kill}, \cmd{pkill}, dan \cmd{killall};
    \item mengelola job di foreground dan background menggunakan \cmd{fg}, \cmd{bg}, \cmd{jobs}, dan \cmd{nohup}.
\end{enumerate}

\begin{notebox}
Seluruh praktek pada bab ini dijalankan pada direktori \dir{~/lab-os/chapter05-processes/}.
Buat direktori ini sebelum memulai:
\begin{lstlisting}[language=bash]
mkdir -p ~/lab-os/chapter05-processes
cd ~/lab-os/chapter05-processes
\end{lstlisting}
\end{notebox}

% ------------------------------------------------------------------------
\section{Konsep Proses dan Thread}
\label{sec:process-concept}

Ketika kamu menjalankan sebuah program, Linux tidak hanya membacakan instruksi program
tersebut ke prosesor begitu saja. Sistem operasi membuat sebuah \textbf{proses}: wadah
terisolasi yang berisi program yang sedang dieksekusi, beserta semua sumber daya yang
dibutuhkannya. Ruang memori proses ini sepenuhnya terpisah dari proses lain. Jika satu
proses crash, proses lain tidak terpengaruh, seperti ketika satu tab browser menutup
sendiri tanpa menutup tab-tab lainnya.

Setiap proses memiliki sejumlah atribut yang selalu bisa kamu lihat dengan perintah
\cmd{ps}:

\begin{itemize}
    \item \textbf{PID}: nomor unik pengenal proses di seluruh sistem.
    \item \textbf{PPID}: PID dari proses induk yang menciptakannya.
    \item \textbf{UID} dan \textbf{GID}: identitas pengguna dan grup pemilik proses, menentukan hak aksesnya.
    \item \textbf{Status}: kondisi proses saat ini, apakah sedang berjalan, menunggu, atau dijeda.
    \item \textbf{Nice}: nilai prioritas penjadwalan, berkisar dari -20 hingga +19.
\end{itemize}

Proses baru di Linux lahir melalui mekanisme \textbf{fork-exec}. Saat kamu mengetik
perintah di terminal, shell melakukan \texttt{fork()}: menyalin dirinya sendiri menjadi
proses anak. Proses anak itu kemudian melakukan \texttt{exec()}: mengganti isi memorinya
dengan program yang ingin dijalankan. Dua langkah ini memungkinkan shell mengatur
redirection dan variabel lingkungan sebelum program baru mulai berjalan.

Informasi tentang setiap proses yang sedang berjalan tersimpan di direktori virtual
\dir{/proc/}. Setiap proses memiliki subdirektori bernama sesuai PID-nya, misalnya
\dir{/proc/1234/}. Di dalamnya terdapat berkas-berkas virtual yang berisi detail proses
tersebut: baris perintah yang digunakan (\file{cmdline}), variabel lingkungan
(\file{environ}), penggunaan memori (\file{status}), dan deskriptor berkas yang terbuka
(\dir{fd/}).

Di dalam sebuah proses, bisa terdapat beberapa alur eksekusi yang berjalan secara
bersamaan. Alur eksekusi ini disebut \textbf{thread}. Tidak seperti proses, thread dalam
satu proses berbagi ruang memori dan sumber daya yang sama. Membuat thread jauh lebih
murah secara komputasi daripada membuat proses baru, karena tidak perlu mengalokasikan
ruang memori terpisah. Server web seperti Apache atau Nginx memanfaatkan ini untuk
melayani ratusan permintaan secara bersamaan dalam satu proses.

\subsection*{Praktek 5.1 Amati Hierarki Proses dengan ps dan pstree}

Tujuan praktek ini adalah melihat hierarki proses secara langsung dan memahami hubungan
induk-anak antarproses menggunakan \cmd{ps} dan \cmd{pstree}.

\begin{enumerate}
    \item Tampilkan semua proses yang berjalan di sistem:
\begin{lstlisting}[language=bash, caption={Menampilkan semua proses dengan ps aux}]
ps aux
\end{lstlisting}
    \textit{Amati:} Kolom \texttt{PID} adalah nomor unik proses. Kolom \texttt{PPID} menunjukkan siapa induknya. Kolom \texttt{STAT} menunjukkan status proses. Kolom \texttt{COMMAND} menampilkan perintah yang menjalankan proses tersebut.

    \item Lihat PID shell yang sedang aktif dan informasi proses induknya:
\begin{lstlisting}[language=bash, caption={Melihat PID shell dan hierarki induknya}]
echo "PID shell ini: $$"
ps -p $$ -f
\end{lstlisting}
    \textit{Amati:} Variabel \texttt{\$\$} berisi PID shell yang sedang berjalan. Kolom \texttt{PPID} pada output menunjukkan PID proses yang membuka shell ini, biasanya program terminal atau SSH daemon.

    \item Tampilkan hierarki proses secara visual menggunakan \cmd{pstree}:
\begin{lstlisting}[language=bash, caption={Melihat pohon proses dengan pstree}]
pstree -p | head -30
\end{lstlisting}

    \item Baca informasi proses langsung dari filesystem virtual:
\begin{lstlisting}[language=bash, caption={Membaca informasi proses dari /proc}]
cat /proc/$$/cmdline | tr '\0' ' '
echo ""
ls /proc/$$/fd | head -10
\end{lstlisting}
    \textit{Amati:} Berkas \file{/proc/PID/cmdline} menyimpan perintah yang menjalankan proses tersebut. Karakter pemisah di dalamnya adalah null byte (\texttt{\textbackslash0}), sehingga perlu dikonversi dengan \cmd{tr} agar terbaca.
\end{enumerate}

\begin{challengebox}
Tampilkan seluruh pohon proses yang berawal dari PID 1 (init atau systemd) menggunakan
kombinasi \cmd{pstree -p 1} dan pipa ke \cmd{grep} (yang sudah kamu pelajari di Bab~3).
Gunakan grep untuk menyaring dan menampilkan hanya baris yang mengandung nama proses
\texttt{sshd}, \texttt{bash}, atau \texttt{cron}. Kemudian gunakan \cmd{ps -p} dengan
salah satu PID yang kamu temukan untuk menampilkan detail proses tersebut.
\end{challengebox}

% ------------------------------------------------------------------------
\section{Siklus Hidup Proses}
\label{sec:process-lifecycle}

Proses tidak langsung berjalan dan berhenti. Ia melewati serangkaian kondisi (\textit{state})
selama hidupnya. Kondisi ini menggambarkan apa yang sedang dilakukan proses tersebut pada
suatu saat. Kamu bisa melihat kondisi setiap proses pada kolom \texttt{STAT} di output
\cmd{ps aux}.

Berikut kondisi-kondisi proses yang akan sering kamu temui. Ringkasannya dapat dilihat
pada \cref{tab:process-states}:

\begin{table}[htbp]
    \centering
    \caption{Kondisi proses di Linux dan artinya}
    \label{tab:process-states}
    \begin{tabular}{lll}
        \hline
        \textbf{Kode} & \textbf{Nama} & \textbf{Keterangan} \\
        \hline
        R & Running/Runnable  & Sedang dieksekusi CPU atau antri untuk dieksekusi \\
        S & Sleeping          & Menunggu event; bisa diinterupsi sinyal \\
        D & Uninterruptible   & Menunggu I/O disk; tidak bisa diinterupsi sinyal \\
        T & Stopped           & Dijeda oleh sinyal SIGSTOP atau debugger \\
        Z & Zombie            & Selesai tapi belum dibersihkan proses induknya \\
        \hline
    \end{tabular}
\end{table}

Kondisi yang paling normal adalah \texttt{S} (sleeping). Sebagian besar proses di sistem
menghabiskan mayoritas waktunya dalam kondisi ini, menunggu input dari pengguna, menunggu
data dari disk, atau menunggu respons dari jaringan. Proses hanya berubah ke kondisi
\texttt{R} ketika ia benar-benar sedang menggunakan prosesor.

Kondisi \texttt{Z} atau zombie perlu perhatian khusus. Proses zombie adalah proses yang
sudah selesai dieksekusi, tetapi entri-nya di tabel proses belum dihapus karena proses
induknya belum memanggil \texttt{wait()} untuk mengambil exit code-nya. Proses zombie
tidak menggunakan memori atau prosesor, tetapi tetap menempati slot di tabel proses.
Jika proses induk mati sebelum membersihkan zombie, proses zombie tersebut diadopsi oleh
proses dengan PID 1 (systemd) yang kemudian membersihkannya.

Kondisi \texttt{D} menandakan proses sedang menunggu operasi I/O yang tidak bisa
diinterupsi. Banyaknya proses dalam kondisi ini biasanya menandakan ada masalah dengan
disk atau sistem penyimpanan yang lambat merespons.

\subsection*{Praktek 5.2 Buat dan Amati Proses di Background}

Tujuan praktek ini adalah mengamati kondisi proses secara langsung dan memahami
bagaimana kondisi proses berubah berdasarkan apa yang sedang dilakukannya.

\begin{enumerate}
    \item Jalankan beberapa proses di latar belakang dan amati kondisinya:
\begin{lstlisting}[language=bash, caption={Menjalankan proses di background}]
cd ~/lab-os/chapter05-processes
sleep 120 &
sleep 240 &
echo "PID proses sleep pertama: $!"
ps aux | grep sleep | grep -v grep
\end{lstlisting}
    \textit{Amati:} Variabel \texttt{\$!} menyimpan PID dari proses background terakhir yang baru dijalankan. Perhatikan kolom \texttt{STAT} pada output \cmd{ps}: kondisi \texttt{S} menandakan proses sedang tidur menunggu timer selesai.

    \item Jeda sebuah proses dengan SIGSTOP dan amati perubahan kondisinya:
\begin{lstlisting}[language=bash, caption={Menjeda proses dan mengamati perubahan status}]
PID_SLEEP=$(pgrep -n sleep)
kill -SIGSTOP $PID_SLEEP
ps aux | grep sleep | grep -v grep
\end{lstlisting}
    \textit{Amati:} Setelah menerima SIGSTOP, kondisi proses berubah dari \texttt{S} menjadi \texttt{T} (stopped). Proses tidak lagi aktif menunggu apapun, ia hanya berhenti sepenuhnya.

    \item Lanjutkan proses yang dijeda dan amati perubahannya kembali:
\begin{lstlisting}[language=bash, caption={Melanjutkan proses yang dijeda}]
kill -SIGCONT $PID_SLEEP
ps aux | grep sleep | grep -v grep
\end{lstlisting}

    \item Bersihkan proses-proses percobaan:
\begin{lstlisting}[language=bash, caption={Membersihkan proses percobaan}]
pkill sleep 2>/dev/null
echo "Proses sleep dibersihkan"
\end{lstlisting}
\end{enumerate}

\begin{challengebox}
Jalankan tiga proses \cmd{sleep} di background dengan durasi berbeda (60, 120, dan 180
detik). Gunakan \cmd{ps aux} dan \cmd{grep} (dari Bab~2) untuk menampilkan hanya baris
yang berisi kata \texttt{sleep}. Kemudian gunakan \cmd{ps -o pid,ppid,stat,ni,cmd}
untuk menampilkan kolom-kolom spesifik saja bagi ketiga proses tersebut. Apa nilai
\texttt{PPID} ketiga proses itu? Apa artinya?
\end{challengebox}

% ------------------------------------------------------------------------
\section{Prioritas Proses}
\label{sec:process-scheduling}

Linux menjalankan puluhan bahkan ratusan proses secara bersamaan, padahal jumlah inti
prosesor terbatas. Kernel menggunakan komponen yang disebut \textbf{scheduler} untuk
memutuskan proses mana yang mendapat giliran waktu prosesor. Scheduler modern Linux
menggunakan algoritma yang berupaya berlaku adil, tetapi kamu bisa memengaruhi keputusan
scheduler dengan mengatur prioritas proses.

Prioritas proses dikendalikan oleh nilai \textbf{nice}. Nama ini berasal dari gagasan
bahwa proses yang ``baik hati'' rela mengalah dan memberi prioritas lebih tinggi ke proses
lain. Nilai nice berkisar dari \textbf{-20} (prioritas tertinggi, paling agresif) hingga
\textbf{+19} (prioritas terendah, paling mengalah). Nilai default adalah \textbf{0}.

Bayangkan ini seperti antrian di warung makan. Nilai nice yang rendah seperti -10 artinya
proses itu bisa menyela antrian dan dilayani lebih cepat. Nilai nice yang tinggi seperti
+15 artinya proses itu sabar mengalah dan baru dilayani setelah yang lain terlayani.
Secara default, semua proses berdiri di antrian dengan sabar di posisi yang sama (nice 0).

Ada aturan penting tentang siapa yang bisa mengatur nilai nice. Pengguna biasa hanya bisa
\textit{menaikkan} nilai nice (menurunkan prioritas) untuk proses miliknya sendiri. Mengatur
nilai nice menjadi negatif, yang meningkatkan prioritas, memerlukan hak akses \texttt{root}.
Pembatasan ini mencegah pengguna biasa dari membuat proses yang memonopoli prosesor.

Perintah \cmd{nice} digunakan saat \textit{memulai} proses baru dengan nilai nice tertentu.
Perintah \cmd{renice} digunakan untuk mengubah nilai nice proses yang \textit{sudah berjalan}.

\begin{lstlisting}[language=bash, caption={Sintaks nice dan renice}]
# Menjalankan perintah dengan nice +10
nice -n 10 perintah argumen

# Mengubah nice proses yang sudah berjalan
renice -n 15 -p 1234

# Melihat nilai nice semua proses (kolom NI)
ps -eo pid,ni,cmd | head -20
\end{lstlisting}

Selain prioritas prosesor, ada juga \cmd{ionice} untuk mengatur prioritas akses I/O disk.
Perintah ini berguna ketika kamu menjalankan operasi yang intensif membaca atau menulis
disk, seperti backup besar, dan tidak ingin operasi itu memperlambat akses disk proses lain.

\subsection*{Praktek 5.3 Ubah Prioritas Proses yang Berjalan}

Tujuan praktek ini adalah menggunakan \cmd{nice} dan \cmd{renice} untuk mengatur
prioritas proses dan memverifikasi perubahannya dengan \cmd{ps}.

\begin{enumerate}
    \item Jalankan tiga proses dengan nilai nice yang berbeda:
\begin{lstlisting}[language=bash, caption={Menjalankan proses dengan nice berbeda}]
cd ~/lab-os/chapter05-processes
nice -n 0  sleep 300 &
nice -n 10 sleep 300 &
nice -n 19 sleep 300 &
echo "Tiga proses sleep dijalankan"
\end{lstlisting}

    \item Tampilkan nilai nice ketiga proses dan amati perbedaannya:
\begin{lstlisting}[language=bash, caption={Melihat nilai nice setiap proses}]
ps -eo pid,ni,cmd | grep sleep | grep -v grep
\end{lstlisting}
    \textit{Amati:} Kolom \texttt{NI} menampilkan nilai nice setiap proses. Ketiga proses menjalankan perintah yang sama (\cmd{sleep}), tetapi memiliki prioritas penjadwalan yang berbeda.

    \item Ubah nilai nice proses pertama menjadi 15 menggunakan \cmd{renice}:
\begin{lstlisting}[language=bash, caption={Mengubah nice proses yang sedang berjalan}]
PID_PERTAMA=$(pgrep -n sleep)
renice -n 15 -p $PID_PERTAMA
ps -p $PID_PERTAMA -o pid,ni,cmd
\end{lstlisting}

    \item Coba atur nilai nice negatif sebagai pengguna biasa:
\begin{lstlisting}[language=bash, caption={Mencoba nice negatif tanpa root}]
renice -n -5 -p $PID_PERTAMA
echo "Exit code: $?"
\end{lstlisting}
    \textit{Amati:} Perintah di atas akan gagal dengan pesan error ``Permission denied'' karena mengatur nilai nice negatif memerlukan hak akses root.

    \item Bersihkan semua proses percobaan:
\begin{lstlisting}[language=bash, caption={Membersihkan proses percobaan}]
pkill sleep
\end{lstlisting}
\end{enumerate}

\begin{challengebox}
Jalankan tiga proses \cmd{sleep 500} di background, masing-masing dengan nilai nice 0,
10, dan 18. Kemudian jalankan \cmd{top} dalam mode batch selama satu detik menggunakan
\cmd{top -bn1} dan arahkan outputnya ke berkas \file{top-snapshot.txt}. Gunakan \cmd{grep}
untuk menyaring hanya baris yang mengandung kata \texttt{sleep} dari berkas tersebut.
Bandingkan kolom \texttt{NI} dan \texttt{PR} (priority) ketiga proses: apakah nilai
\texttt{PR} berbanding lurus dengan nilai nice? Bersihkan semua proses setelah selesai.
\end{challengebox}

% ------------------------------------------------------------------------
\section{Process Signal}
\label{sec:process-signals}

Bagaimana cara berkomunikasi dengan proses yang sedang berjalan? Kamu tidak bisa mengetik
langsung ke dalamnya jika proses itu berjalan di latar belakang tanpa antarmuka interaktif.
Jawabannya adalah \textbf{\term{signal}}. \term{Signal} adalah notifikasi singkat yang dikirim ke sebuah
proses untuk memintanya melakukan sesuatu: berhenti, dijeda, memuat ulang konfigurasi,
atau tindakan lainnya.

Setiap \term{signal} punya nama dan nomor. Kamu bisa menggunakan keduanya secara bergantian.
Signal yang paling sering digunakan diringkas pada \cref{tab:important-signals}:

\begin{table}[htbp]
    \centering
    \caption{Signal penting di Linux}
    \label{tab:important-signals}
    \begin{tabular}{llll}
        \hline
        \textbf{Nama} & \textbf{No.} & \textbf{Shortcut} & \textbf{Keterangan} \\
        \hline
        SIGHUP  & 1  &                    & Hangup: minta reload konfigurasi \\
        SIGINT  & 2  & \keystroke{Ctrl+C} & Interupsi dari keyboard \\
        SIGKILL & 9  &                    & Paksa berhenti, tidak bisa ditangkap \\
        SIGTERM & 15 &                    & Minta berhenti dengan bersih (default) \\
        SIGSTOP & 19 & \keystroke{Ctrl+Z} & Jeda proses, tidak bisa ditangkap \\
        SIGCONT & 18 &                    & Lanjutkan proses yang dijeda \\
        \hline
    \end{tabular}
\end{table}

Perbedaan terpenting adalah antara \textbf{SIGTERM} dan \textbf{SIGKILL}. SIGTERM adalah
permintaan sopan kepada proses untuk berhenti. Proses bisa menangkap \term{signal} ini, menyimpan
data yang belum tersimpan, menutup koneksi yang terbuka, lalu berhenti dengan bersih.
SIGKILL adalah perintah paksa dari kernel: proses tidak bisa menolak atau menangkap \term{signal}
ini. Kernel langsung mengakhiri proses tanpa memberi kesempatan cleanup. Selalu coba
SIGTERM terlebih dahulu; gunakan SIGKILL hanya jika proses tidak merespons SIGTERM setelah
beberapa detik.

Perintah \cmd{kill} menerima PID sebagai argumen. Meski namanya ``kill'', fungsinya adalah
mengirim sinyal apa pun, bukan hanya sinyal untuk mematikan proses. Perintah \cmd{pkill}
lebih praktis karena menerima nama proses, bukan PID. Perintah \cmd{killall} mengirim
sinyal ke semua proses yang namanya cocok secara persis. Perintah \cmd{pgrep} berguna
untuk mencari PID berdasarkan nama sebelum mengirim sinyal dengan \cmd{kill}.

\begin{lstlisting}[language=bash, caption={Cara-cara mengirim sinyal}]
# Kirim SIGTERM (default) berdasarkan PID
kill 1234

# Kirim SIGKILL berdasarkan PID
kill -9 1234

# Kirim sinyal berdasarkan nama proses
pkill nama-proses
pkill -9 nama-proses

# Cari PID berdasarkan nama proses
pgrep sleep
\end{lstlisting}

\subsection*{Praktek 5.4 Kirim Sinyal ke Proses}

Tujuan praktek ini adalah mengirim berbagai jenis sinyal ke proses dan mengamati
efeknya pada kondisi proses.

\begin{enumerate}
    \item Buat tiga proses percobaan:
\begin{lstlisting}[language=bash, caption={Membuat proses percobaan}]
cd ~/lab-os/chapter05-processes
sleep 500 &
sleep 600 &
sleep 700 &
ps aux | grep sleep | grep -v grep
\end{lstlisting}

    \item Kirim SIGTERM ke proses pertama dan verifikasi ia berhenti:
\begin{lstlisting}[language=bash, caption={Menghentikan proses dengan SIGTERM}]
PID_500=$(pgrep -f "sleep 500")
kill $PID_500
ps aux | grep sleep | grep -v grep
\end{lstlisting}
    \textit{Amati:} Proses \cmd{sleep 500} tidak lagi muncul di daftar. SIGTERM berhasil menghentikannya.

    \item Jeda proses kedua dengan SIGSTOP dan amati perubahan kondisinya:
\begin{lstlisting}[language=bash, caption={Menjeda proses dengan SIGSTOP}]
PID_600=$(pgrep -f "sleep 600")
kill -SIGSTOP $PID_600
ps aux | grep sleep | grep -v grep
\end{lstlisting}
    \textit{Amati:} Kondisi proses \cmd{sleep 600} berubah dari \texttt{S} menjadi \texttt{T}. Proses tidak berjalan maupun menunggu event, ia benar-benar dijeda di tempat.

    \item Lanjutkan proses yang dijeda dengan SIGCONT:
\begin{lstlisting}[language=bash, caption={Melanjutkan proses dengan SIGCONT}]
kill -SIGCONT $PID_600
ps aux | grep sleep | grep -v grep
\end{lstlisting}

    \item Hentikan semua proses sleep sekaligus:
\begin{lstlisting}[language=bash, caption={Menghentikan semua proses sleep sekaligus}]
pkill sleep
ps aux | grep sleep | grep -v grep
\end{lstlisting}
\end{enumerate}

\begin{challengebox}
Tulis skrip sederhana bernama \file{proses-tahan.sh} yang menangkap sinyal SIGTERM dan
menampilkan pesan sebelum keluar. Isi skrip tersebut adalah:

\begin{lstlisting}[language=bash, caption={Skrip yang menangkap SIGTERM}]
#!/bin/bash
trap 'echo "Menerima SIGTERM, keluar dengan bersih..."; exit 0' SIGTERM
echo "Proses berjalan dengan PID: $$"
while true; do
    sleep 1
done
\end{lstlisting}

Beri izin eksekusi menggunakan \cmd{chmod u+x} (dari Bab~4), jalankan di background,
lalu kirim SIGTERM ke PID-nya. Amati pesan yang muncul. Ini adalah gambaran awal tentang
penanganan sinyal dalam skrip yang akan kamu pelajari lebih lanjut di Bab~7.
\end{challengebox}

% ------------------------------------------------------------------------
\section{Manajemen Job}
\label{sec:job-management}

Dalam konteks shell, \textbf{job} adalah proses yang dikelola langsung oleh shell tersebut.
Setiap job bisa berjalan di \textbf{foreground} (menguasai terminal, kamu harus menunggu
selesai sebelum bisa mengetik perintah lain) atau di \textbf{background} (berjalan di
belakang layar, terminal tetap bebas digunakan).

Tambahkan tanda \texttt{\&} di akhir perintah untuk menjalankannya langsung di background.
Jika sebuah proses sudah terlanjur berjalan di foreground dan ternyata butuh waktu lama,
kamu bisa menekan \keystroke{Ctrl+Z} untuk menjeda dan memindahkannya ke background
menggunakan \cmd{bg}.

Berikut perintah-perintah manajemen job yang perlu kamu kuasai. Ringkasannya dapat
dilihat pada \cref{tab:job-management-commands}:

\begin{table}[htbp]
    \centering
    \caption{Perintah manajemen job di shell}
    \label{tab:job-management-commands}
    \begin{tabular}{ll}
        \hline
        \textbf{Perintah} & \textbf{Fungsi} \\
        \hline
        \texttt{perintah \&}    & Jalankan di background sejak awal \\
        \cmd{jobs}              & Tampilkan semua job aktif di shell ini \\
        \cmd{fg \%N}            & Bawa job nomor N ke foreground \\
        \cmd{bg \%N}            & Lanjutkan job nomor N di background \\
        \keystroke{Ctrl+Z}      & Jeda dan pindahkan job foreground ke background \\
        \keystroke{Ctrl+C}      & Hentikan job foreground \\
        \cmd{kill \%N}          & Hentikan job nomor N \\
        \hline
    \end{tabular}
\end{table}

Ada satu masalah dengan job biasa: ketika kamu menutup terminal (atau logout dari SSH),
semua job yang dijalankan dari terminal itu menerima sinyal SIGHUP dan berhenti. Ini
masalah ketika kamu ingin menjalankan proses panjang seperti pembuatan laporan atau
pemrosesan data besar, lalu menutup terminal dan pergi.

Solusinya adalah \cmd{nohup}. Perintah ini menjalankan program dengan mengabaikan
SIGHUP. Ditambah dengan \texttt{\&} untuk background, proses tersebut akan tetap berjalan
meski kamu menutup terminal. Perintah \cmd{disown} melakukan hal serupa untuk proses yang
sudah terlanjur berjalan: ia melepas proses dari daftar job shell sehingga proses tidak
menerima SIGHUP saat shell ditutup.

\begin{lstlisting}[language=bash, caption={Menjalankan proses tahan logout}]
# Jalankan program dengan nohup, output ke nohup.out
nohup perintah &

# Jalankan dengan output khusus
nohup perintah > output.log 2>&1 &

# Lepas job yang sudah ada dari shell
disown %1
\end{lstlisting}

\subsection*{Praktek 5.5 Kelola Beberapa Job Sekaligus}

Tujuan praktek ini adalah mengelola job foreground dan background, serta menjalankan
proses yang tetap berjalan meski terminal ditutup menggunakan \cmd{nohup}.

\begin{enumerate}
    \item Jalankan tiga job di background dan tampilkan daftarnya:
\begin{lstlisting}[language=bash, caption={Menjalankan tiga job di background}]
cd ~/lab-os/chapter05-processes
sleep 300 &
sleep 400 &
sleep 500 &
jobs
\end{lstlisting}
    \textit{Amati:} Setiap job mendapat nomor (ditampilkan dalam kurung siku). Tanda \texttt{+} menandakan job terakhir yang aktif, tanda \texttt{-} menandakan job sebelumnya.

    \item Bawa job pertama ke foreground:
\begin{lstlisting}[language=bash, caption={Membawa job ke foreground}]
fg %1
\end{lstlisting}

    \item Jeda job yang sedang di foreground dengan \keystroke{Ctrl+Z}, lalu kembalikan ke background:
\begin{lstlisting}[language=bash]
# Tekan Ctrl+Z di sini
bg %1
jobs
\end{lstlisting}
    \textit{Amati:} Setelah \cmd{bg \%1}, job kembali berjalan di background. Status di output \cmd{jobs} berubah dari \texttt{Stopped} menjadi \texttt{Running}.

    \item Jalankan proses panjang dengan \cmd{nohup} agar tetap berjalan meski terminal ditutup:
\begin{lstlisting}[language=bash, caption={Menjalankan proses tahan logout}]
nohup sleep 3600 > nohup-test.log 2>&1 &
echo "PID nohup: $!"
jobs
\end{lstlisting}

    \item Lepas proses nohup dari shell dan verifikasi ia tetap berjalan:
\begin{lstlisting}[language=bash, caption={Melepas proses dari shell dengan disown}]
disown %$(jobs | tail -1 | grep -oP '^\[\K[0-9]+')
jobs
ps aux | grep "sleep 3600" | grep -v grep
\end{lstlisting}
    \textit{Amati:} Setelah \cmd{disown}, proses tidak lagi muncul di output \cmd{jobs}, tetapi masih ada di output \cmd{ps aux}. Proses ini tidak akan berhenti meski terminal ditutup.

    \item Bersihkan semua proses percobaan:
\begin{lstlisting}[language=bash, caption={Membersihkan semua proses percobaan}]
kill %1 %2 %3 2>/dev/null
pkill -f "sleep 3600" 2>/dev/null
jobs
\end{lstlisting}
\end{enumerate}

\begin{challengebox}
Jalankan perintah berikut di background menggunakan \cmd{nohup}:
\begin{lstlisting}[language=bash]
nohup bash -c 'for i in $(seq 1 60); do echo "$(date): iterasi $i"; sleep 2; done' \
    > ~/lab-os/chapter05-processes/proses-panjang.log 2>&1 &
\end{lstlisting}
Gunakan \cmd{tail -f} untuk memantau berkas log secara langsung sambil proses berjalan.
Tekan \keystroke{Ctrl+C} untuk menghentikan \cmd{tail -f} (bukan prosesnya). Kemudian
gunakan \cmd{ps aux} dan \cmd{grep} untuk memverifikasi proses masih berjalan. Setelah
selesai, hentikan prosesnya dengan \cmd{pkill -f proses-panjang}.
\end{challengebox}

% ------------------------------------------------------------------------
\section{Rangkuman}
\label{sec:process-summary}

Dalam bab ini, kamu telah mempelajari:

\begin{itemize}[noitemsep,topsep=2pt]
    \item \textbf{Proses dan thread}: Proses adalah instance program yang berjalan dengan ruang memori terpisah dan PID unik. Thread berbagi memori dalam satu proses dan lebih ringan secara komputasi. Informasi proses tersedia di \dir{/proc/PID/}.
    \item \textbf{Siklus hidup proses}: Proses melewati kondisi R (running), S (sleeping), D (uninterruptible), T (stopped), dan Z (zombie). Proses baru lahir melalui mekanisme fork-exec. Kondisi proses bisa dilihat di kolom \texttt{STAT} pada output \cmd{ps aux}.
    \item \textbf{Prioritas proses}: Nilai nice berkisar dari -20 (prioritas tertinggi) hingga +19 (terendah). Perintah \cmd{nice} mengatur prioritas saat memulai proses, \cmd{renice} mengubah prioritas proses yang sudah berjalan. Nilai negatif memerlukan akses root.
    \item \textbf{Sinyal proses}: Sinyal adalah notifikasi ke proses. SIGTERM (15) meminta proses berhenti dengan bersih, SIGKILL (9) menghentikan secara paksa. Gunakan \cmd{kill} berdasarkan PID, \cmd{pkill} berdasarkan nama proses.
    \item \textbf{Manajemen job}: Tambahkan \texttt{\&} untuk background, \cmd{fg}/\cmd{bg} untuk berpindah antara foreground dan background. Perintah \cmd{nohup} dan \cmd{disown} memastikan proses tetap berjalan meski terminal ditutup.
\end{itemize}

% ------------------------------------------------------------------------
\section{Latihan}
\label{sec:latihan-proses}

\textbf{Instruksi Umum}: Kerjakan seluruh latihan secara mandiri. Catat langkah penting, simpan tangkapan layar bila diperlukan, lalu rangkum hasilnya sebagai dokumentasi pribadi.

\subsubsection*{Latihan 5.1 Eksplorasi Hierarki Proses}
Jalankan \cmd{ps aux --forest} dan temukan proses dengan PID 1. Jawab pertanyaan berikut
dengan bukti tangkapan layar: (1) Apa nama proses PID 1 dan apa fungsinya di sistem
Linux modern? (2) Berapa total proses yang dimiliki oleh user root dan berapa yang
dimiliki oleh usermu? (3) Gunakan \cmd{pstree -p} dan temukan PID dari shell yang sedang
aktif. Proses apa yang menjadi induknya?


\subsubsection*{Latihan 5.2 Simulasi Manajemen Job}
Jalankan lima perintah \cmd{sleep} dengan durasi 100, 200, 300, 400, dan 500 detik di
background. Dokumentasikan langkah-langkah berikut: (1) Tampilkan daftar semua job
dengan \cmd{jobs}. (2) Bawa job ketiga ke foreground, jeda dengan \keystroke{Ctrl+Z},
lalu kembalikan ke background. (3) Hentikan job pertama dan kedua menggunakan \cmd{kill}.
(4) Tampilkan kembali daftar job yang tersisa. Sertakan tangkapan layar setiap langkah.


\subsubsection*{Latihan 5.3 Prioritas dan Sinyal}
Lakukan eksperimen berikut dan dokumentasikan hasilnya: (1) Jalankan dua proses
\cmd{sleep} dengan nice +5 dan nice +15. Verifikasi nilai NI keduanya. (2) Ubah nice
proses pertama menjadi +10 menggunakan \cmd{renice}. Setelah diubah, proses mana yang
sekarang lebih diprioritaskan? (3) Kirim SIGSTOP ke salah satu proses, verifikasi kondisi
\texttt{T}-nya di output \cmd{ps}, lalu kirim SIGCONT. (4) Akhiri semua proses percobaan
dengan \cmd{pkill sleep}. Sertakan semua perintah dan tangkapan layar.


% ------------------------------------------------------------------------
\section{Referensi}
\label{sec:referensi-proses}

\begin{itemize}
    \item Kerrisk, M. \textit{The Linux Programming Interface} (1st ed.). No Starch Press, 2010.
    \item Ward, B. \textit{How Linux Works} (3rd ed.). No Starch Press, 2021.
    \item Barrett, D., Silverman, R., Byrnes, R. \textit{Linux Pocket Guide} (3rd ed.). O'Reilly Media, 2016.
    \item Halaman manual: \texttt{man ps}, \texttt{man top}, \texttt{man 7 signal}, \texttt{man nice}, \texttt{man renice}, \texttt{man kill}, \texttt{man pkill}, \texttt{man nohup}, \texttt{man pgrep}.
\end{itemize}
